\scanpage{127}
\uline{NB} Finalement, on n'a pas eu besoin du tout de faire intervenir $\rho$ dans la vérification, c'était une complication inutile

\begin{center}\rule{3cm}{0.4pt}\end{center}

Je vais reprendre la théorie précédente. Soient $H_{\deleted{\unsure{k}}},H'$ deux fibrés vectoriels de rang 2 \unsure{sur|\deleted{sur}} $k$, munis d'une \deleted{forme $q$} structure de $k$-algèbre (néc. comm.) d'où \deleted{forme}\inserted{sur} $H$ et $H'$ formes \uline{norme} et trace
\diagram{\begin{tikzcd}[ampersand replacement=\&,column sep=small,row sep=small]
{\deleted{\Delta}N_H:H}\arrow[r]\& k\&{\deleted{\Delta}N_{H'}:H'}\arrow[r,dash]\&k\\ {T_H:H}\arrow[r]\&k\&{T_{H'}:H'}\arrow[r,dash]\&k\\ {\varepsilon_H:\deleted{\unsure{1_H}}k}\arrow[r]\&H\&{\varepsilon_{H'}:k'}\arrow[r,dash]\&H'
\end{tikzcd}}
% slanted
\margin{$\varepsilon_H(\lambda)=\lambda1_H$ \quad $\varepsilon_{H'}(\lambda)=\unsure{\rho_0|\deleted{\rho_0}}\lambda.1_{H'}$} \deleted{l'essentiel de ce qui va suivre, n'aura pas besoin de la structure multiplicative, mais seulement des formes précédentes (voire, seulement de $N$ et $T_H$)} La formule
% slanted
\margin{\uline{NB} La forme $T_H$ se déduit de la donnée de $1_H\in H$ (i.e. $\varepsilon_H$) et de $\Delta_H$, comme $T_H(u)=\varphi_H(u,1)$ et le fait que $\Delta_H(1)=1$ (d'où $T(1)=2$) implique que les f\textsuperscript{es} $\varphi_H(1,u)=\psi_H(1,u,)\quad\forall u$}
\[
\underbrace{\varphi_{\deleted{\unsure{N|\Delta}}H}}_{=\varphi_{\Delta_H}}(u,v)=\operatorname{Tr}(u)\operatorname{Tr}(v)-\underbrace{\operatorname{Tr}(uv)}_{\psi_H(u,v)}
\]
montre que la forme $\operatorname{Tr}(uv)$ sur $H$ s'exprime par la connaissance de $N_H$ et $T_H$ comme
\[
\operatorname{Tr}_{H}(uv)=T_H(u)T_H(v)-\varphi_{\deleted{\Delta}H}(u,v)
\]
i.e.
\[
\varphi_H(u,v)+\psi_H(u,v)=T_H(u)T_H(v)
\]
% archive sheet number: 126
